% Conclusion. Draft prose --- framing stable; keep claims scoped to null-controlled
% evidence. Fill the final effect size at revision.

Sampling-based hallucination detectors treat semantic uncertainty as a property of the
question, invariant to surface form. We tested that assumption adversarially, on the
detector itself. Our evidence indicates it can be made to fail under a
meaning-preserving paraphrase---verified equivalent by bidirectional NLI, leaving the
model's correctness unchanged---which, reported net of a benign-paraphrase noise floor,
inflates the score of a \emph{correct} answer into an apparent false alarm, the
under-explored direction we foreground. The mirror \emph{hide} direction, which
suppresses a wrong answer's score, the attack is built to exercise but we do not yet
report under the same null control. We were careful to make this claim only where the
evidence supports it: targets are chosen without reference to the detector's own scores,
success requires the hallucination status to survive the paraphrase, and effect sizes
are reported net of a benign-paraphrase noise floor rather than as the raw output of an
optimiser that maximises a noisy estimate.

Two boundaries remain, and we state them as the work to do rather than as caveats to
wave away. First, our clusterer and our equivalence gate are the same NLI model, so
whether the vulnerability is a property of the \emph{paradigm} or of one model's
self-inconsistency turns on an independent clustering oracle. We build one: the cheap
candidates (exact match, sentence embeddings) fail on the adversarial case, but a
victim- and NLI-independent LLM-judge, validated on exactly that case ($0.92$ vs.\
$0.51$), clears it---and it is under this judge that the confirmatory large-$n$ run,
now under way, adjudicates whether the effect survives the independent clusterer, a
verdict we draw from that run rather than from the machinery-validation sample. Second,
the meaning-preserving guarantee rests on an automatic gate; a human equivalence audit
is needed before it is fully established. A
simple reformulation-averaging defense raises the attacker's cost but does not close
the gap, suggesting that robustness to paraphrase must be designed into these detectors
rather than assumed. Taken together, the results argue that ``semantic'' uncertainty
inherits the blind spots of whatever model defines its notion of meaning, and that
detectors built on a learned equivalence relation should be evaluated as the
adversarial targets they now are.
